\section{Introduction}
\label{sec:intro}

The modern landscape of deep learning is defined by an explosion of specialized models. Rather than training colossal, monolithic networks from scratch for every novel task, the community has turned to Parameter-Efficient Fine-Tuning (PEFT) techniques, such as Low-Rank Adaptation (LoRA) \cite{hu2021lora}, to specialize shared pre-trained backbones \cite{devlin2018bert, dosovitskiy2020image, touvron2023llama}. However, this success has introduced a formidable deployment bottleneck: how do we serve a multitude of specialized task-specific adapters simultaneously to a stream of heterogeneous, noisy, and unpredictable real-world requests on edge-compute hardware?

To address this, two main paradigms have emerged: static model merging and dynamic routing. Static merging methods, such as Task Arithmetic \cite{ilharco2022editing}, TIES-Merging \cite{yadav2023ties}, and DARE \cite{yu2023dare}, compress multiple specialized expert adapters into a single static set of weights. While parameter-free and efficient, these methods are fundamentally limited by representation overlap and structural conflict; they yield a "jack-of-all-trades, master-of-none" network that suffers from mutual interference and severe task-accuracy degradation under diverse task workloads. 
Dynamic routing frameworks, such as SABLE and SPS-ZCA, attempt to resolve this by dynamically routing activations to specific experts or blending adapter parameters on-the-fly based on intermediate representations.

\begin{figure*}[t]
    \centering
    \begin{subfigure}[b]{0.48\textwidth}
        \centering
        \includegraphics[width=\textwidth]{noise_sensitivity.png}
        \caption{Domain Noise Scaling Stress Test}
        \label{fig:noise_sensitivity}
    \end{subfigure}
    \hfill
    \begin{subfigure}[b]{0.48\textwidth}
        \centering
        \includegraphics[width=\textwidth]{batch_size_heterogeneity.png}
        \caption{Deployment Stream Batch Sweep}
        \label{fig:batch_size}
    \end{subfigure}
    \caption{Key experimental results of ESM-LVC in the 14-layer Isolating Coordinate Sandbox (ICS). (a) Under severe scaling domain noise (Scale 2.5), ESM-LVC preserves an outstanding Joint Mean accuracy of 65.37\%, outperforming SPS-ZCA by +2.63\% absolute. (b) While parametric Linear Router suffers from severe heterogeneity collapse, ESM-LVC maintains flatline, robust performance from micro-batches of size $B=1$ up to $B=512$.}
    \label{fig:top_level_results}
\end{figure*}

Despite their improvements, these existing dynamic ensembling methods are built upon a fundamentally reductionist and static assumption: they treat task-specific experts as isolated, independent entities that exist in a vacuum, ignoring the semantic relationships and structural co-existence of these experts within the shared parameter space. These methods formulate routing as a static feedforward mapping or projection (e.g., via simple cosine similarities or linear routing heads). Consequently, when deployed in non-ideal real-world environments, they are highly sensitive to out-of-domain noise: a small amount of input corruption or domain shift blurs the routing coordinates, causing massive activation misrouting, parameter-space interference, and performance collapse.

In this work, we propose a radical paradigm shift. We ask: \emph{What if we view the mixture of specialized experts not as a static, isolated set of channels, but as a living, self-organizing symbiotic ecosystem?} 

We argue that task-specific experts fine-tuned from a shared backbone naturally possess underlying semantic affinities and conflicts. Instead of using a static, feedforward routing head, we can model the activation levels of these experts as interacting biological populations. Under this ecological perspective, similar or complementary tasks should exhibit \textbf{mutualism (cooperation)}, where they reinforce each other's activation channels to boost generalization, while dissimilar or conflicting tasks should exhibit \textbf{competitive exclusion (mutual suppression)}, where they actively suppress each other to prevent representation collapse and shield task specialization.

To turn this vision into a concrete, computationally efficient framework, we introduce \textbf{Evolutionary Symbiotic Merging via Lotka-Volterra Cooperation (ESM-LVC)}. ESM-LVC implements classical Lotka-Volterra competition-cooperation differential equations directly into the activation space of a neural network. It consists of three tightly coupled, training-free, and parameter-free components:
\begin{itemize}
    \item \textbf{Lotka-Volterra Activation Dynamics (LVAD):} A non-linear dynamical system that governs how the ensembling coefficients (representing expert populations) evolve over a virtual localized time scale $\tau$ inside the forward pass.
    \item \textbf{Symbiotic Interaction Tensor (SIT):} A pre-computed semantic interaction matrix $\Gamma$ that defines the cooperative (mutualistic) and competitive (exclusionary) relationships between experts based on their intermediate centroid alignments.
    \item \textbf{Discrete Euler Symbiosis Solver (DESS):} An ultra-lightweight discrete solver that integrates these differential equations inside a single forward pass with negligible computational overhead, requiring only a few element-wise vector-matrix operations.
\end{itemize}

To validate the mathematical and dynamic properties of our framework, we conduct a calibrated proof-of-concept numerical simulation study in a 192-dimensional Isolating Coordinate Sandbox (ICS) analytically designed to emulate the activation-subspace and power-law performance characteristics of a 14-layer Vision Transformer (ViT-Tiny) backbone serving four simulated vision task adapters (MNIST, Fashion-MNIST, CIFAR-10, SVHN). We explicitly frame this evaluation as a numerical simulation study rather than a physical deep learning deployment, establishing a highly controlled environment to isolate the mathematical properties of our Lotka-Volterra dynamics under noise and task overlap. Under this calibrated simulation, our findings reveal several profound properties of ESM-LVC:
\begin{enumerate}
    \item \textbf{State-of-the-Art Performance:} Under standard settings, ESM-LVC achieves \textbf{75.12\%} Joint Mean accuracy, surpassing existing non-parametric baselines such as SABLE (74.13\%) and SPS-ZCA (74.31\%).
    \item \textbf{Self-Regulating Noise Resilience:} Under extreme scaling domain noise (Scale 2.5), while other methods suffer from severe routing degradation, ESM-LVC maintains an outstanding \textbf{65.37\%} accuracy, outperforming SOTA SPS-ZCA by \textbf{+2.63\%} absolute. The non-linear Lotka-Volterra dynamics act as a natural high-pass filter, suppressing weak, noise-driven out-of-domain activations.
    \item \textbf{Inherent Immunity to Heterogeneity Collapse:} Under deployment stream sweeps from $B=1$ to $B=512$, ESM-LVC exhibits flatline, robust performance, whereas parametric routing heads collapse completely.
\end{enumerate}

By bringing non-linear dynamical systems and symbiotic ecology directly into the forward pass of neural networks, ESM-LVC demonstrates that biological principles can inspire resilient, robust, and training-free model serving. While our physical verification (Section~\ref{subsec:physical_verification}) bridges the simulation gap by testing the solver offline on real activations from a pre-trained Vision Transformer, we explicitly highlight that end-to-end active physical adapter blending remains our immediate next empirical milestone.
